% ID: SEL-D0081-OB-E04
% ESTADO: cotejada. Segunda pasada visual y borrador compilado revisados.
% Fuente DOC: Selectividad/sel_2012_matematicas_aplicadas/3 MATEMÁT. APLIC. CC. SS/EXAMEN 2-ANDALUCÍA.doc
% SHA-256 original: bcec57ae525fa5e23ce74631274a2abda0b5b16ac47ce6108ba8b5572d4aa634
% Página convertida: 2; opción B; ejercicio 4.

\begin{enumerate}[label=\alph*)]
\item \textbf{(1 punto)} En una ciudad viven 400 hombres y 320 mujeres y se quiere seleccionar una muestra de tamaño 54 utilizando muestreo estratificado por sexos, con afijación proporcional, \textquestiondown{}cuál sería la composición de la muestra?
\item \textbf{(1.5 puntos)} A partir de una población de elementos 1, 2, 3, 4 se seleccionan, mediante muestreo aleatorio simple, todas las muestras de tamaño 2.

Escriba dichas muestras y calcule la varianza de las medias muestrales.
\end{enumerate}
